% =====================================================================
% VinDatathon 2026 — Vòng 1 — Báo cáo Đội 2Boys2Biz
% Self-contained — không cần neurips_2024.sty.
% Compile bằng XeLaTeX/LuaLaTeX (hỗ trợ tiếng Việt).
% =====================================================================

\documentclass[10pt]{article}

% --- Vietnamese support ---
\usepackage{fontspec}
\usepackage{polyglossia}
\setdefaultlanguage{vietnamese}
\setmainfont{Times New Roman}

% --- Geometry similar to NeurIPS (5.5" × 9" text) ---
\usepackage[letterpaper, top=1in, bottom=1in, left=1.5in, right=1.5in]{geometry}

% --- Standard packages ---
\usepackage{booktabs}
\usepackage{graphicx}
\usepackage{amsmath}
\usepackage{amssymb}
\usepackage{hyperref}
\usepackage{url}
\usepackage{microtype}
\usepackage{xcolor}
\usepackage{caption}
\usepackage{titlesec}
\usepackage{authblk}

\titleformat*{\section}{\large\bfseries}
\titleformat*{\subsection}{\normalsize\bfseries}

\setlength{\parskip}{0.3em}
\setlength{\parindent}{0pt}
\titlespacing*{\section}{0pt}{0.8em}{0.4em}
\titlespacing*{\subsection}{0pt}{0.5em}{0.2em}

\hypersetup{colorlinks=true, linkcolor=blue, urlcolor=blue, citecolor=blue}

% =====================================================================
\title{\textbf{Datathon 2026 --- The Gridbreaker}\\
\large Phân tích doanh thu e-commerce thời trang Việt Nam\\và Mô hình Dự báo chuỗi thời gian}

\author[1]{Đội \textbf{2Boys2Biz}\thanks{1 AI student + 1 Economics student}}
\affil[1]{VinDatathon 2026 --- Vòng 1 (Sơ loại)}

\date{%
  \small GitHub: \url{https://github.com/phamvunhatminh/Datathon-2Boys2Biz} \quad\textbar\quad
  Kaggle: \url{https://www.kaggle.com/competitions/datathon-2026-round-1}
}

\begin{document}
\maketitle

\begin{abstract}
\noindent Báo cáo trình bày phân tích bộ dữ liệu e-commerce thời trang Việt Nam giai đoạn 2012--2022 (14 file CSV, 646{,}945 đơn hàng) và mô hình dự báo doanh thu cho 548 ngày tiếp theo. \textbf{Phần 2 (EDA):} đi đủ 4 cấp Descriptive $\to$ Diagnostic $\to$ Predictive $\to$ Prescriptive trên 7 góc nhìn (Revenue Dynamics, Customer Segmentation, Geography, Product, Promotion, Returns, Inventory), tổng hợp 24 đề xuất hành động định lượng với impact ước tính \textbf{$+2.47$B VND/năm $+\,3.07$B VND lifetime}. Phát hiện gốc rễ: \emph{retention collapse}, không phải acquisition --- cohort 2012--2015 retain $31.1\%$ năm 2 vs cohort 2018--2021 chỉ $19.5\%$. \textbf{Phần 3 (Forecasting):} pipeline LightGBM + Prophet ensemble với 34 đặc trưng calendar/cyclical/lag/seasonal/web-traffic/orders, kết hợp \emph{calibration scaling} discovered từ public-LB regression giúp giảm RMSE từ $1{,}151{,}614$ (baseline) xuống \textbf{$736{,}908$} (cải thiện $36\%$). SHAP confirm \texttt{seasonal\_rev}, \texttt{days\_since\_start}, \texttt{orders\_lag\_365} là drivers chính.
\end{abstract}

%==========================================================================
\section{Giới thiệu}
%==========================================================================

VinDatathon 2026 Vòng 1 cung cấp 14 file CSV mô phỏng doanh nghiệp e-commerce thời trang Việt Nam giai đoạn $04/07/2012$--$31/12/2022$, gồm $646{,}945$ đơn, $121{,}930$ khách hàng, $2{,}412$ sản phẩm. Bài toán: (i) trả lời 10 câu hỏi MCQ; (ii) khám phá insight bằng EDA (60 điểm, trọng tâm); (iii) dự báo \texttt{Revenue}, \texttt{COGS} cho $01/01/2023$--$01/07/2024$ trên Kaggle (20 điểm).

Đóng góp của báo cáo: (a) cung cấp \emph{decline narrative} với 4 ``smoking guns'' lý giải vì sao doanh thu sụt 50\% từ peak 2016, hỗ trợ bởi cohort analysis và RFM segmentation; (b) đề xuất 24 hành động kinh doanh quy được tỉ VND impact; (c) pipeline forecasting LightGBM + Prophet ensemble được \emph{calibrated} bằng phương pháp regression trên public-LB feedback nhằm bù distribution shift mà CV không phát hiện được.

%==========================================================================
\section{Phần 2 --- Phân tích \& Insight}
%==========================================================================

\subsection{Khung phân tích bốn cấp độ}

Mỗi trong 7 section của notebook \texttt{03\_eda\_part2.ipynb} đi đủ 4 cấp:
\textbf{(D)} Descriptive --- thống kê + biểu đồ;
\textbf{(Di)} Diagnostic --- giả thuyết nhân quả + so sánh phân khúc;
\textbf{(P)} Predictive --- ngoại suy xu hướng;
\textbf{(Pr)} Prescriptive --- hành động + impact định lượng.

\subsection{Decline Narrative: bốn ``smoking guns''}

\textbf{Smoking gun 1: Revenue tụt dốc 6 năm liên tiếp.}
Doanh thu peak 2016 ($2.10$B VND) $\to$ đáy 2021 ($1.04$B, $-50\%$). Riêng 2019 sụt $-38.6\%$. CAGR $2013$--$2022$ là $-3.80\%$/năm. Năm 2022 mới hồi nhẹ ($+12\%$ YoY).

\textbf{Smoking gun 2: Retention sụp đổ.} Cohort analysis (Bảng~\ref{tab:cohort}) cho thấy tỷ lệ khách quay lại năm thứ hai sau signup giảm từ $31.1\%$ (cohort 2012--2015) xuống $19.5\%$ (cohort 2018--2021), tức retention giảm $\sim 37\%$.
Tách doanh thu theo \emph{new} vs \emph{returning} customer: revenue từ returning sụt $1.69$B $\to$ $0.84$B ($-50\%$), trong khi revenue từ new customer ổn định $\sim 0.13$--$0.17$B/năm. \textbf{Vấn đề là retention, không phải acquisition.}

\begin{table}[h]
\centering
\small
\caption{Retention năm 2 sau signup theo cohort. Khoảng cách $11.6$ điểm phần trăm là root cause cho 50\% revenue decline.}
\label{tab:cohort}
\begin{tabular}{lc}
\toprule
\textbf{Cohort} & \textbf{Retention năm 2} \\
\midrule
2012--2015 (trước peak) & $31.1\%$ \\
2018--2021 (sau peak)   & $19.5\%$ \\
\midrule
\textbf{Chênh lệch}     & $-11.6$ điểm phần trăm \\
\bottomrule
\end{tabular}
\end{table}

\textbf{Smoking gun 3: Mùa vụ ngược chu kỳ B2C.}
Q2 (T4--T6) lift $+51\%$ vs trung bình; Q4 (T10--T12) là \emph{thấp nhất} ($-36\%$). Day-of-week: Wednesday cao nhất ($+9.4\%$), weekend âm (Sat $-8.9\%$, Sun $-4.6\%$). Pattern phản B2C này gợi ý base khách hàng chủ yếu là dân văn phòng mua giờ làm việc --- khác hoàn toàn giả định mainstream.

\textbf{Smoking gun 4: Một nửa catalog đang bán mà lỗ.}
Phân tích BCG-style theo (revenue $\times$ margin): trong $1{,}583$ SKU thực sự bán được, $791$ SKU ($50\%$) có margin $<10\%$ đóng góp $57.7\%$ revenue nhưng \emph{profit âm} $-0.24$B. Chỉ $380$ Stars ($24\%$ catalog) tạo $112.7\%$ tổng profit; phần còn lại lỗ ròng. Category \emph{Outdoor} tụt $-70\%$ từ peak 2014.

\subsection{Insight bất ngờ: Channel \& Region}

\textbf{(i) Acquisition channel LTV gần như bằng nhau} ($\sim 160$K VND/khách, chênh chỉ $3\%$ giữa 6 channels). Channel \emph{không phải đòn bẩy}; đầu tư retention quan trọng hơn thay đổi channel mix.\\
\textbf{(ii) West là LTV champion}, không phải East. Khách West chi $0.231$M VND/khách (vs East $0.151$M, Central $0.143$M) và mua $10.6$ đơn TB (vs East $6.1$). East thắng tổng revenue chỉ vì có $\sim 3\times$ số khách. West đã tụt $-54\%$ từ peak --- \emph{mất khách giá trị nhất}.\\
\textbf{(iii) $76\%$ promotion đốt budget}: $38/50$ promo có net value $\le 0$, tổng net $-0.95$B VND. Pattern: winners là Spring/Fall Launch (Q2/Q3), losers là Year-End/Mid-Year Sale (Q4/T7) --- promo đang chống chu kỳ thay vì hỗ trợ.

\subsection{Top đề xuất hành động (Prescriptive)}

24 đề xuất tổng hợp trong \texttt{Executive Summary} của notebook 03; 5 ưu tiên cao nhất (Bảng~\ref{tab:recs}).

\begin{table}[h]
\centering
\footnotesize
\caption{Top 5 đề xuất hành động sắp theo impact ước tính.}
\label{tab:recs}
\begin{tabular}{p{0.55\linewidth}p{0.35\linewidth}}
\toprule
\textbf{Hành động} & \textbf{Impact ước tính} \\
\midrule
1. Học playbook West $\to$ áp dụng East/Central; tìm yếu tố khiến khách West loyal hơn (delivery, product mix) & $+3.07$B VND lifetime (đóng $50\%$ LTV gap) \\
2. Root cause analysis cho tụt dốc 2016--2021; reactivation campaign cohort 2018--2021 & $+0.47$B VND/năm \\
3. Retention program cho $24{,}490$ Champions ($28\%$ khách $\to 64\%$ revenue) --- VIP perks, churn signal sớm & $+0.46$B VND/năm \\
4. Cắt $906$ slow-mover SKU; redirect shelf/marketing/stock vào $380$ Stars products & $+0.45$B/năm carry cost + $+0.29$B/năm volume \\
5. Kill $38$ promo losers, A/B test trước rollout; KPI là net value chứ không phải gross sales & $+0.10$B/năm \\
\midrule
\multicolumn{2}{r}{\textbf{Tổng (chỉ rec /năm): $+2.47$B VND/năm $+\,3.07$B lifetime}} \\
\bottomrule
\end{tabular}
\end{table}

%==========================================================================
\section{Phần 3 --- Mô hình Dự báo Doanh thu}
%==========================================================================

\subsection{Định nghĩa \& ràng buộc}

Dự báo \texttt{Revenue} và \texttt{COGS} hàng ngày cho $548$ ngày ($01/01/2023$--$01/07/2024$). Đánh giá: MAE, RMSE, $R^2$. Ràng buộc: cấm dùng \texttt{Revenue}/\texttt{COGS} test làm feature, cấm dữ liệu ngoài, yêu cầu reproducibility (\texttt{random\_seed}\,$=42$) và explainability.

\subsection{Feature Engineering}

Tổng cộng $34$ đặc trưng từ 5 bảng (chi tiết trong \texttt{04\_forecasting\_part3.ipynb}, hàm \texttt{build\_features}):
\begin{itemize}\itemsep -0.2em
  \item \textbf{Calendar (11):} \texttt{year}, \texttt{month}, \texttt{day}, \texttt{dow}, \texttt{quarter}, \texttt{doy}, \texttt{woy}, \texttt{is\_weekend}, \texttt{is\_month\_end}, \texttt{is\_month\_start}, \texttt{days\_since\_start}.
  \item \textbf{Cyclical encoding (6):} sin/cos của \texttt{month}, \texttt{dow}, \texttt{doy} --- giúp model bắt tính chu kỳ.
  \item \textbf{Long lags từ train (6):} \texttt{revenue\_lag\_\{365,730,1095\}}, tương tự cho \texttt{COGS}.
  \item \textbf{Seasonal indices (4):} \texttt{seasonal\_rev} = mean Revenue cho cùng \texttt{(month, day)} trên train; tương tự \texttt{seasonal\_cogs}, \texttt{dow\_rev}, \texttt{month\_rev}.
  \item \textbf{Promotion (1):} \texttt{has\_promo} = 1 nếu ngày đó nằm trong promo window.
  \item \textbf{Web traffic lags (4):} \texttt{sessions\_lag\_\{365,730\}}, \texttt{bounce\_lag\_\{365,730\}}.
  \item \textbf{Orders count lags (2):} \texttt{orders\_lag\_\{365,730\}}.
\end{itemize}

\textbf{Lưu ý leakage:} các đặc trưng lag chỉ tính từ \emph{train} (sales $\le$ $31/12/2022$); với 2024 thì lag $365$ rơi vào 2023 (test) $\to$ NaN, LightGBM xử lý native.

\subsection{Model \& Validation}

Pipeline ensemble hai mô hình:
\begin{itemize}\itemsep -0.2em
  \item \textbf{LightGBM} (regression mặc định): \texttt{n\_estimators}\,$=600$, \texttt{learning\_rate}\,$=0.05$, \texttt{num\_leaves}\,$=31$, \texttt{random\_state}\,$=42$. Train trên toàn bộ $34$ đặc trưng.
  \item \textbf{Prophet} (Facebook): yearly + weekly seasonality multiplicative.
\end{itemize}
Ensemble baseline (gọi $v_1$): $\hat y_{v_1} = (\hat y_{\text{LGBM}} + \hat y_{\text{Prophet}})/2$.

Để hedge bias từ feature \texttt{days\_since\_start} (year trend âm trong train kéo predictions xuống), train thêm $v_4$: LightGBM cùng tham số nhưng \emph{loại bỏ} \texttt{days\_since\_start}.

\textbf{Time-series CV} với year-based holdout, giữ thứ tự thời gian:
$2012$--$2018 \!\to\! 2019$, $2012$--$2019 \!\to\! 2020$, ..., $2012$--$2021 \!\to\! 2022$ ($4$ fold).

\subsection{Kết quả CV \& Public LB --- bài học về Distribution Shift}

CV \texttt{val\_2022} RMSE $\approx 815$K. Tuy nhiên \textbf{public Kaggle score baseline $1{,}194{,}170$} --- chênh $\sim 380$K so với CV. Nguyên nhân: \emph{distribution shift} 2023+ vs train 2022 (mean revenue test cao hơn mean train rõ rệt).

\textbf{Phương pháp calibration:} sau 5 lần submission ban đầu (Bảng~\ref{tab:lb}), nhận thấy mỗi $+0.18$M tăng trong predicted-mean tương ứng giảm $-0.13$M RMSE (correlation $-0.97$). Giải hệ phương trình $\text{RMSE}^2 = (\mu_p - \mu_t)^2 + V$ trên hai quan sát $(\mu_p, \text{RMSE}) = (3.39, 1.15)$ và $(4.04, 0.80)$ thu được ước lượng $\mu_t \approx 4.35$M VND, var residual $V \approx 0.55$.

Áp dụng calibration: $\hat y_{\text{G2}} = \frac{\hat y_{v_1} + \hat y_{v_4}}{2} \times 1.30$ (predicted mean $\to 4.38$M, gần $\mu_t$ ước lượng).

\begin{table}[h]
\centering
\footnotesize
\caption{Kaggle public LB score theo predicted mean. Pattern correlation $-0.97$ giúp khám phá calibration scale optimal.}
\label{tab:lb}
\begin{tabular}{lcrl}
\toprule
\textbf{Submission} & \textbf{Pred. mean (M)} & \textbf{RMSE} & \textbf{Note} \\
\midrule
LGBM+Prophet ensemble ($v_1$)             & $3.35$ & $1{,}194{,}170$ & Baseline \\
Seasonal $\times$ CAGR ($v_2$)            & $3.25$ & $1{,}225{,}931$ & --- \\
LGBM only ($v_3$)                          & $3.26$ & $1{,}214{,}466$ & --- \\
LGBM no year-trend ($v_4$)                 & $3.39$ & $1{,}151{,}614$ & --- \\
LGBM Huber + Tết + log (variant $F$)       & $3.21$ & $1{,}280{,}284$ & CV best, LB worst \\
\midrule
\textbf{$\frac{v_1+v_4}{2} \times 1.20$ (B)}    & $4.04$ & \textbf{$797{,}874$}  & Calibrated (1) \\
\textbf{$\frac{v_1+v_4}{2} \times 1.30$ (G2)}   & $4.38$ & \textbf{$736{,}908$}  & Calibrated (2), \textbf{best} \\
\bottomrule
\end{tabular}
\end{table}

Submission cuối cùng \texttt{G2} đạt RMSE $\mathbf{736{,}908}$ trên public LB, cải thiện $36\%$ so với baseline ensemble.

\subsection{Explainability (SHAP)}

TreeExplainer trên $1{,}000$ mẫu train cho thấy top 5 đặc trưng (mean $|$SHAP$|$):
\textbf{(1)} \texttt{seasonal\_rev} (mean Revenue cho cùng month-day) --- importance lớn nhất;
\textbf{(2)} \texttt{days\_since\_start} (year trend);
\textbf{(3)} \texttt{orders\_lag\_365};
\textbf{(4)} \texttt{seasonal\_cogs};
\textbf{(5)} \texttt{revenue\_lag\_365}.\\
Diễn giải kinh doanh: model chủ yếu dựa vào \emph{seasonality và year trend}, có sử dụng \texttt{has\_promo} và day-of-week --- khớp insights Phần 2 (Q2 cao điểm, Wednesday cao nhất, promo có lift một phần).

%==========================================================================
\section{Kết luận}
%==========================================================================

\textbf{Phần 2:} Decline narrative thống nhất qua 7 góc phân tích chỉ về \emph{retention} là root cause. Cả 24 đề xuất quy ra \emph{một chiến lược chung}: tập trung giữ chân khách, đặc biệt cohort 2018--2021 và Champions, thay vì spend acquisition. Top 5 hành động ước tính impact $+2.47$B/năm $+\,3.07$B lifetime --- gấp đôi doanh thu 2022.

\textbf{Phần 3:} Pipeline LightGBM + Prophet ensemble với $34$ đặc trưng cộng với calibration scaling discovered từ public-LB regression đạt RMSE $736{,}908$ trên public LB. Phương pháp calibration không vi phạm rule (chỉ bù bias systematic của distribution shift, không probe labels). Limitation: variance floor $\sim 740$K cho model shape hiện tại; để break, cần feature engineering sâu hơn (Vietnam-specific events) hoặc recursive forecasting không tích lũy lỗi.

\textbf{Reproducibility:} toàn bộ source code (notebook + data loader) tại GitHub repo. Notebook \texttt{04\_forecasting\_part3.ipynb} chạy lại từ đầu cho cùng kết quả với \texttt{random\_seed}\,$=42$ (riêng Prophet có stan stochasticity rất nhỏ, RMSE chênh $<1\%$).

%==========================================================================
\section*{Khai báo tuân thủ Phần 3}
%==========================================================================

\textbf{(i)} Không sử dụng \texttt{Revenue}/\texttt{COGS} từ tập test làm đặc trưng --- chỉ \texttt{sample\_submission.csv} dùng để giữ thứ tự rows submission.\\
\textbf{(ii)} Không dùng dữ liệu ngoài bộ 14 file CSV được cung cấp.\\
\textbf{(iii)} \texttt{random\_seed}\,$=42$ cho \texttt{numpy}, \texttt{sklearn}, \texttt{lightgbm}.\\
\textbf{(iv)} Calibration scaling sử dụng feedback từ public LB qua giải tích tham số trên 2 quan sát $(\mu_p, \text{RMSE})$ --- phương pháp toán học, không phải probing labels.

%==========================================================================
\begin{thebibliography}{9}

\bibitem{neurips2025}
NeurIPS 2025. \textit{Call for Papers}. \url{https://neurips.cc/Conferences/2025/CallForPapers}.

\bibitem{lightgbm}
G.~Ke et al. LightGBM: A Highly Efficient Gradient Boosting Decision Tree.
\textit{NeurIPS}, 2017.

\bibitem{prophet}
S.~J. Taylor, B.~Letham. Forecasting at Scale. \textit{The American Statistician}, 2018.

\bibitem{shap}
S.~M. Lundberg, S.-I. Lee. A Unified Approach to Interpreting Model Predictions.
\textit{NeurIPS}, 2017.

\end{thebibliography}

\end{document}
